%% Experiences section.

\section{\faBriefcase$\;$\textbf{Expériences professionnelles}}
\cventry{Août 2019 $\rightharpoonup$\\Maintenant}{\href{https://www.metroscope.tech}{Metroscope}}{FullStack Java/React Developer}{Paris (17ème)}{}{%
  \ifthenelse{\equal{\displayImageExperiences}{yes}}{
    \vspace{-0.2cm}
    \begin{wrapfigure}{R}{0.2\textwidth}%
      \vspace{-0.3cm}
      \centering%
      \href{https://www.metroscope.tech}{\includegraphics[width=0.18\textwidth]{./img/metroscope}}%
      \vspace{+0.5cm}
    \end{wrapfigure}%
  }{}
Developpeur fullstack de l'équipe Dev de Metroscope (2 back, 2 front, 2 fullstack). Division du backend (Java \& SpringBoot) en micro-services. Mutualisation des configurations Gradle via l'écriture de plugins (Groovy), nettoyage des dépendances. Réécriture from scratch du client React en TypeScript, utilisation des Hooks, internationalisation, réflexion autour de l'UX/UI.\\Stack technique: React, TypeScript, Java, Gradle, Docker.}

\vspace{\cvExperienceSpace}

\cventry{Mars 2018 $\rightharpoonup$\\Juillet 2019}{\href{https://www.wemanity.com}{Wemanity}}{Agile Consultant Developer R\&D}{Paris (10ème)}{}{%
  \ifthenelse{\equal{\displayImageExperiences}{yes}}{
    \vspace{-0.2cm}
    \begin{wrapfigure}{R}{0.2\textwidth}%
      \vspace{-0.5cm}
      \centering%
      \href{https://www.wemanity.com}{\includegraphics[width=0.18\textwidth]{./img/wemanity}}%
    \end{wrapfigure}%
  }{}
  R\&D Lead Developer au sein du Lab Innovation \textit{Innov8} de la \href{https://www.societegenerale.fr/}{Société Générale}.\newline
  \faCaretRight\hspace{0.5mm}Idéation, développement et maintenance d'applications autour des robots \href{https://www.softbankrobotics.com/emea/en/robots/pepper}{\textit{Pepper}} de \href{https://www.softbankrobotics.com/emea/en}{Softbank Robotics} pour des évènements, de l'animation et cas d'usage internes.\newline
  \faCaretRight\hspace{0.5mm}Développement de prototypes d’applications de data visualisation mêlant modélisation en 3D/VR et graphes de données métier complexes.\newline
  \faCaretRight\hspace{0.5mm}Étude et développement de services cognitifs (\href{https://fr.wikipedia.org/wiki/Reconnaissance_optique_de_caractères}{OCR}, \href{https://fr.wikipedia.org/wiki/Compréhension_du_langage_naturel}{NLU} et \href{https://fr.wikipedia.org/wiki/Reconnaissance_automatique_de_la_parole}{STT}) à l'aide de solutions open source.\\Stack technique: ReactJS, ES6, Loopback, Python, Docker, OpenShift.}

\vspace{\cvExperienceSpace}

\cventry{Octobre 2016 $\rightharpoonup$\\Février 2018}{\href{https://www.dhatim.com}{Dhatim}}{Développeur Java R\&D}{Paris (12ème)}{}{%
  \ifthenelse{\equal{\displayImageExperiences}{yes}}{
    \vspace{-0.2cm}
    \begin{wrapfigure}{R}{0.2\textwidth}%
      \vspace{-0.3cm}
      \centering%
      \href{https://www.dhatim.com}{\includegraphics[width=0.18\textwidth]{./img/dhatim}}%
    \end{wrapfigure}%
  }{}
  Développeur backend rattaché à l'équipe de \href{https://pay.conciliator.ai/}{\textit{Conciliator For Pay}}, solution en SaaS de contrôle de DSN (Déclaration Sociale Nominative) pour les entreprises. Équipe composée de 5 personnes (2 front, 2 back, 1 PO) organisée en méthode Scrum.\\Stack technique: Java 8 Core, Dropwizard, PostgreSQL, Docker, Git, Jenkins.}

\vspace{\cvExperienceSpace}

\cventry{Mars 2015 $\rightharpoonup$\\Octobre 2016}{\href{https://www.visian.fr}{Visian}}{Lead Developer \& Responsable Fablab}{Nanterre (92)}{}{%
  \ifthenelse{\equal{\displayImageExperiences}{yes}}{
    \vspace{-0.2cm}
    \begin{wrapfigure}{R}{0.2\textwidth}%
      \vspace{-0.4cm}
      \centering%
      \href{https://www.visian.fr}{\includegraphics[width=0.18\textwidth]{./img/visian}}%
    \end{wrapfigure}%
  }{}
  Lead Dev \& Lead Tech chez Visian, la Business Unit R\&D Internet of Things intégrée au groupe \href{https://www.neurones-it.com}{Neurones IT}. En charge de l'équipe des développeurs (hardware \& software). Réalisation et développement de PoC mêlant IoT (beacons, smartwatch, lunettes connectées, Raspberry Pi, Arduino, NFC.) et cloud (Microsoft Azure, IBM Bluemix, AWS IoT) pour la remontée et le traitement des données.\\Stack technique: C, Java 8 Core, Android, Swift, Python, Git.}

\vspace{\cvExperienceSpace}

\cventry{Septembre 2014 $\rightharpoonup$\\Mars 2015}{\href{https://www.groupe-sii.com/fr/}{Groupe SII}}{Assistant Responsable Innovation}{Paris (75)}{}{%
  \ifthenelse{\equal{\displayImageExperiences}{yes}}{
    \vspace{-0.2cm}
    \begin{wrapfigure}{R}{0.2\textwidth}%
      \vspace{-0.8cm}
      \centering%
      \href{https://www.groupe-sii.com/fr/}{\includegraphics[width=0.18\textwidth]{./img/sii}}%
    \end{wrapfigure}%
  }{}
  Étude et exploration des drones open source, prototypage d'un mini drone quadrirotor avec l'intégration d'une caméra et d'un radio-transmetteur vidéo (\href{https://www.bitcraze.se/}{Crazyflie project}). Mise en place d’un prototype d'objet connecté à partir d’un \href{https://raspberrypicomputer.com/}{Raspberry Pi} d'interrupteurs, diodes et capteurs (température, pression, GPS).\\Environnement technique: C, Java 8 ME for IoT, XMPP.}

\vspace{\cvExperienceSpace}

\cventry{Avril 2013 $\rightharpoonup$\\Août 2013}{\href{https://www.thalesgroup.fr/}{Thales Services}}{Stagiaire Administrateur Système}{Ramonville (31)}{}{%
  \ifthenelse{\equal{\displayImageExperiences}{yes}}{
    \vspace{-0.2cm}
    \begin{wrapfigure}{R}{0.2\textwidth}%
      \vspace{-0.3cm}
      \centering%
      \href{https://www.thalesgroup.fr/}{\includegraphics[width=0.18\textwidth]{./img/thales}}%
    \end{wrapfigure}%
  }{}
  Responsable de l'étude, de l'installation et du déploiement d'\href{https://www.openstack.org}{OpenStack} dans un environnement de tests. Rédaction du dossier d’architecture matérielle, réalisation des scripts de déploiement automatique, tests et validations.\\Environnement technique : OpenStack, Mirantis, RHEL, Bash scripting.}

\vspace{\cvExperienceSpace}

\cventry{Octobre 2012 $\rightharpoonup$\\Avril 2013}{\href{https://www.novacom-services.com/}{Novacom Services}}{Stagiaire Développeur}{Ramonville (31)}{}{%
  \ifthenelse{\equal{\displayImageExperiences}{yes}}{
    \vspace{-0.2cm}
    \begin{wrapfigure}{R}{0.2\textwidth}%
      \vspace{-0.8cm}
      \centering%
      \href{https://www.novacom-services.com/}{\includegraphics[width=0.18\textwidth]{./img/novacom}}%
    \end{wrapfigure}%
  }{}
  Développement d'un jeu de données et d'un ensemble de tests unitaires et fonctionnels sur la couche service des applications, permettant l’automatisation des validations logicielles et ainsi garantir la non régression pendant les développements.\\Environnement technique : Java 8, JUnit, Jenkins, Git.}

\vspace{\cvSectionSpace}
